\documentclass[11pt]{article}
\newcommand{\videoid}{jBuvlWhNWkg}
% Preámbulo compartido de las notas del curso de topología algebraica.
% Cada nota empieza con \documentclass[11pt]{article}, define \videoid{...}
% (el identificador del video de YouTube) y luego hace % Preámbulo compartido de las notas del curso de topología algebraica.
% Cada nota empieza con \documentclass[11pt]{article}, define \videoid{...}
% (el identificador del video de YouTube) y luego hace % Preámbulo compartido de las notas del curso de topología algebraica.
% Cada nota empieza con \documentclass[11pt]{article}, define \videoid{...}
% (el identificador del video de YouTube) y luego hace \input{preambulo}.
% Compilar con LuaLaTeX:  lualatex nombre-de-la-nota.tex
\usepackage[spanish,provide=*]{babel}
\usepackage[T1]{fontenc}
\usepackage{lmodern}
\usepackage[a4paper,margin=2.6cm]{geometry}
\usepackage{amsmath,amssymb,amsthm}
\usepackage{xcolor}
\usepackage{tikz}
\usetikzlibrary{arrows.meta,babel}
\usepackage[colorlinks=true,urlcolor=blue!50!black,linkcolor=blue!50!black]{hyperref}

\hyphenpenalty=10000 \exhyphenpenalty=10000 \emergencystretch=3em

% Colores de las figuras (los mismos que usa la página)
\definecolor{extremo}{RGB}{47,107,79}
\definecolor{azul}{RGB}{44,74,154}

% Definiciones numeradas por unidad
\theoremstyle{definition}
\newtheorem{defn}{Definición}
% Numeración por unidad: \unidad{u}{n} (después de \input{preambulo}) hace que las
% definiciones de la nota se llamen u.n, u.(n+1), ...; n es el primer número
% libre de esa unidad (ver la tabla de CONVENCIONES.md).
\newcommand{\unidad}[2]{\renewcommand{\thedefn}{#1.\arabic{defn}}\setcounter{defn}{\numexpr#2-1\relax}}

% Observaciones y erratas del video. Usar SOLO si hay algo que señalar:
%   \begin{erratas}
%     \errata{203}{3:23}{Se dice ..., pero debe decir ...}
%   \end{erratas}
% El primer argumento son los segundos (enlaza al minuto exacto del video),
% el segundo es la marca visible, el tercero es el texto.
\newenvironment{erratas}{\section*{Observaciones y erratas del video}\begin{itemize}}{\end{itemize}}
\newcommand{\errata}[3]{%
  \edef\errurl{https://www.youtube.com/watch?v=\videoid\string&t=#1}%
  \item[\textbf{\expandafter\href\expandafter{\errurl}{#2}}] #3}
.
% Compilar con LuaLaTeX:  lualatex nombre-de-la-nota.tex
\usepackage[spanish,provide=*]{babel}
\usepackage[T1]{fontenc}
\usepackage{lmodern}
\usepackage[a4paper,margin=2.6cm]{geometry}
\usepackage{amsmath,amssymb,amsthm}
\usepackage{xcolor}
\usepackage{tikz}
\usetikzlibrary{arrows.meta,babel}
\usepackage[colorlinks=true,urlcolor=blue!50!black,linkcolor=blue!50!black]{hyperref}

\hyphenpenalty=10000 \exhyphenpenalty=10000 \emergencystretch=3em

% Colores de las figuras (los mismos que usa la página)
\definecolor{extremo}{RGB}{47,107,79}
\definecolor{azul}{RGB}{44,74,154}

% Definiciones numeradas por unidad
\theoremstyle{definition}
\newtheorem{defn}{Definición}
% Numeración por unidad: \unidad{u}{n} (después de % Preámbulo compartido de las notas del curso de topología algebraica.
% Cada nota empieza con \documentclass[11pt]{article}, define \videoid{...}
% (el identificador del video de YouTube) y luego hace \input{preambulo}.
% Compilar con LuaLaTeX:  lualatex nombre-de-la-nota.tex
\usepackage[spanish,provide=*]{babel}
\usepackage[T1]{fontenc}
\usepackage{lmodern}
\usepackage[a4paper,margin=2.6cm]{geometry}
\usepackage{amsmath,amssymb,amsthm}
\usepackage{xcolor}
\usepackage{tikz}
\usetikzlibrary{arrows.meta,babel}
\usepackage[colorlinks=true,urlcolor=blue!50!black,linkcolor=blue!50!black]{hyperref}

\hyphenpenalty=10000 \exhyphenpenalty=10000 \emergencystretch=3em

% Colores de las figuras (los mismos que usa la página)
\definecolor{extremo}{RGB}{47,107,79}
\definecolor{azul}{RGB}{44,74,154}

% Definiciones numeradas por unidad
\theoremstyle{definition}
\newtheorem{defn}{Definición}
% Numeración por unidad: \unidad{u}{n} (después de \input{preambulo}) hace que las
% definiciones de la nota se llamen u.n, u.(n+1), ...; n es el primer número
% libre de esa unidad (ver la tabla de CONVENCIONES.md).
\newcommand{\unidad}[2]{\renewcommand{\thedefn}{#1.\arabic{defn}}\setcounter{defn}{\numexpr#2-1\relax}}

% Observaciones y erratas del video. Usar SOLO si hay algo que señalar:
%   \begin{erratas}
%     \errata{203}{3:23}{Se dice ..., pero debe decir ...}
%   \end{erratas}
% El primer argumento son los segundos (enlaza al minuto exacto del video),
% el segundo es la marca visible, el tercero es el texto.
\newenvironment{erratas}{\section*{Observaciones y erratas del video}\begin{itemize}}{\end{itemize}}
\newcommand{\errata}[3]{%
  \edef\errurl{https://www.youtube.com/watch?v=\videoid\string&t=#1}%
  \item[\textbf{\expandafter\href\expandafter{\errurl}{#2}}] #3}
) hace que las
% definiciones de la nota se llamen u.n, u.(n+1), ...; n es el primer número
% libre de esa unidad (ver la tabla de CONVENCIONES.md).
\newcommand{\unidad}[2]{\renewcommand{\thedefn}{#1.\arabic{defn}}\setcounter{defn}{\numexpr#2-1\relax}}

% Observaciones y erratas del video. Usar SOLO si hay algo que señalar:
%   \begin{erratas}
%     \errata{203}{3:23}{Se dice ..., pero debe decir ...}
%   \end{erratas}
% El primer argumento son los segundos (enlaza al minuto exacto del video),
% el segundo es la marca visible, el tercero es el texto.
\newenvironment{erratas}{\section*{Observaciones y erratas del video}\begin{itemize}}{\end{itemize}}
\newcommand{\errata}[3]{%
  \edef\errurl{https://www.youtube.com/watch?v=\videoid\string&t=#1}%
  \item[\textbf{\expandafter\href\expandafter{\errurl}{#2}}] #3}
.
% Compilar con LuaLaTeX:  lualatex nombre-de-la-nota.tex
\usepackage[spanish,provide=*]{babel}
\usepackage[T1]{fontenc}
\usepackage{lmodern}
\usepackage[a4paper,margin=2.6cm]{geometry}
\usepackage{amsmath,amssymb,amsthm}
\usepackage{xcolor}
\usepackage{tikz}
\usetikzlibrary{arrows.meta,babel}
\usepackage[colorlinks=true,urlcolor=blue!50!black,linkcolor=blue!50!black]{hyperref}

\hyphenpenalty=10000 \exhyphenpenalty=10000 \emergencystretch=3em

% Colores de las figuras (los mismos que usa la página)
\definecolor{extremo}{RGB}{47,107,79}
\definecolor{azul}{RGB}{44,74,154}

% Definiciones numeradas por unidad
\theoremstyle{definition}
\newtheorem{defn}{Definición}
% Numeración por unidad: \unidad{u}{n} (después de % Preámbulo compartido de las notas del curso de topología algebraica.
% Cada nota empieza con \documentclass[11pt]{article}, define \videoid{...}
% (el identificador del video de YouTube) y luego hace % Preámbulo compartido de las notas del curso de topología algebraica.
% Cada nota empieza con \documentclass[11pt]{article}, define \videoid{...}
% (el identificador del video de YouTube) y luego hace \input{preambulo}.
% Compilar con LuaLaTeX:  lualatex nombre-de-la-nota.tex
\usepackage[spanish,provide=*]{babel}
\usepackage[T1]{fontenc}
\usepackage{lmodern}
\usepackage[a4paper,margin=2.6cm]{geometry}
\usepackage{amsmath,amssymb,amsthm}
\usepackage{xcolor}
\usepackage{tikz}
\usetikzlibrary{arrows.meta,babel}
\usepackage[colorlinks=true,urlcolor=blue!50!black,linkcolor=blue!50!black]{hyperref}

\hyphenpenalty=10000 \exhyphenpenalty=10000 \emergencystretch=3em

% Colores de las figuras (los mismos que usa la página)
\definecolor{extremo}{RGB}{47,107,79}
\definecolor{azul}{RGB}{44,74,154}

% Definiciones numeradas por unidad
\theoremstyle{definition}
\newtheorem{defn}{Definición}
% Numeración por unidad: \unidad{u}{n} (después de \input{preambulo}) hace que las
% definiciones de la nota se llamen u.n, u.(n+1), ...; n es el primer número
% libre de esa unidad (ver la tabla de CONVENCIONES.md).
\newcommand{\unidad}[2]{\renewcommand{\thedefn}{#1.\arabic{defn}}\setcounter{defn}{\numexpr#2-1\relax}}

% Observaciones y erratas del video. Usar SOLO si hay algo que señalar:
%   \begin{erratas}
%     \errata{203}{3:23}{Se dice ..., pero debe decir ...}
%   \end{erratas}
% El primer argumento son los segundos (enlaza al minuto exacto del video),
% el segundo es la marca visible, el tercero es el texto.
\newenvironment{erratas}{\section*{Observaciones y erratas del video}\begin{itemize}}{\end{itemize}}
\newcommand{\errata}[3]{%
  \edef\errurl{https://www.youtube.com/watch?v=\videoid\string&t=#1}%
  \item[\textbf{\expandafter\href\expandafter{\errurl}{#2}}] #3}
.
% Compilar con LuaLaTeX:  lualatex nombre-de-la-nota.tex
\usepackage[spanish,provide=*]{babel}
\usepackage[T1]{fontenc}
\usepackage{lmodern}
\usepackage[a4paper,margin=2.6cm]{geometry}
\usepackage{amsmath,amssymb,amsthm}
\usepackage{xcolor}
\usepackage{tikz}
\usetikzlibrary{arrows.meta,babel}
\usepackage[colorlinks=true,urlcolor=blue!50!black,linkcolor=blue!50!black]{hyperref}

\hyphenpenalty=10000 \exhyphenpenalty=10000 \emergencystretch=3em

% Colores de las figuras (los mismos que usa la página)
\definecolor{extremo}{RGB}{47,107,79}
\definecolor{azul}{RGB}{44,74,154}

% Definiciones numeradas por unidad
\theoremstyle{definition}
\newtheorem{defn}{Definición}
% Numeración por unidad: \unidad{u}{n} (después de % Preámbulo compartido de las notas del curso de topología algebraica.
% Cada nota empieza con \documentclass[11pt]{article}, define \videoid{...}
% (el identificador del video de YouTube) y luego hace \input{preambulo}.
% Compilar con LuaLaTeX:  lualatex nombre-de-la-nota.tex
\usepackage[spanish,provide=*]{babel}
\usepackage[T1]{fontenc}
\usepackage{lmodern}
\usepackage[a4paper,margin=2.6cm]{geometry}
\usepackage{amsmath,amssymb,amsthm}
\usepackage{xcolor}
\usepackage{tikz}
\usetikzlibrary{arrows.meta,babel}
\usepackage[colorlinks=true,urlcolor=blue!50!black,linkcolor=blue!50!black]{hyperref}

\hyphenpenalty=10000 \exhyphenpenalty=10000 \emergencystretch=3em

% Colores de las figuras (los mismos que usa la página)
\definecolor{extremo}{RGB}{47,107,79}
\definecolor{azul}{RGB}{44,74,154}

% Definiciones numeradas por unidad
\theoremstyle{definition}
\newtheorem{defn}{Definición}
% Numeración por unidad: \unidad{u}{n} (después de \input{preambulo}) hace que las
% definiciones de la nota se llamen u.n, u.(n+1), ...; n es el primer número
% libre de esa unidad (ver la tabla de CONVENCIONES.md).
\newcommand{\unidad}[2]{\renewcommand{\thedefn}{#1.\arabic{defn}}\setcounter{defn}{\numexpr#2-1\relax}}

% Observaciones y erratas del video. Usar SOLO si hay algo que señalar:
%   \begin{erratas}
%     \errata{203}{3:23}{Se dice ..., pero debe decir ...}
%   \end{erratas}
% El primer argumento son los segundos (enlaza al minuto exacto del video),
% el segundo es la marca visible, el tercero es el texto.
\newenvironment{erratas}{\section*{Observaciones y erratas del video}\begin{itemize}}{\end{itemize}}
\newcommand{\errata}[3]{%
  \edef\errurl{https://www.youtube.com/watch?v=\videoid\string&t=#1}%
  \item[\textbf{\expandafter\href\expandafter{\errurl}{#2}}] #3}
) hace que las
% definiciones de la nota se llamen u.n, u.(n+1), ...; n es el primer número
% libre de esa unidad (ver la tabla de CONVENCIONES.md).
\newcommand{\unidad}[2]{\renewcommand{\thedefn}{#1.\arabic{defn}}\setcounter{defn}{\numexpr#2-1\relax}}

% Observaciones y erratas del video. Usar SOLO si hay algo que señalar:
%   \begin{erratas}
%     \errata{203}{3:23}{Se dice ..., pero debe decir ...}
%   \end{erratas}
% El primer argumento son los segundos (enlaza al minuto exacto del video),
% el segundo es la marca visible, el tercero es el texto.
\newenvironment{erratas}{\section*{Observaciones y erratas del video}\begin{itemize}}{\end{itemize}}
\newcommand{\errata}[3]{%
  \edef\errurl{https://www.youtube.com/watch?v=\videoid\string&t=#1}%
  \item[\textbf{\expandafter\href\expandafter{\errurl}{#2}}] #3}
) hace que las
% definiciones de la nota se llamen u.n, u.(n+1), ...; n es el primer número
% libre de esa unidad (ver la tabla de CONVENCIONES.md).
\newcommand{\unidad}[2]{\renewcommand{\thedefn}{#1.\arabic{defn}}\setcounter{defn}{\numexpr#2-1\relax}}

% Observaciones y erratas del video. Usar SOLO si hay algo que señalar:
%   \begin{erratas}
%     \errata{203}{3:23}{Se dice ..., pero debe decir ...}
%   \end{erratas}
% El primer argumento son los segundos (enlaza al minuto exacto del video),
% el segundo es la marca visible, el tercero es el texto.
\newenvironment{erratas}{\section*{Observaciones y erratas del video}\begin{itemize}}{\end{itemize}}
\newcommand{\errata}[3]{%
  \edef\errurl{https://www.youtube.com/watch?v=\videoid\string&t=#1}%
  \item[\textbf{\expandafter\href\expandafter{\errurl}{#2}}] #3}

\unidad{1}{25} % unidad 1; primer enunciado: 1.25

\title{Homomorfismos inducidos}
\author{Luis Jorge Sánchez Saldaña\\[2pt]\small Facultad de Ciencias, UNAM}
\date{}

\begin{document}
\maketitle

\noindent\small Notas del video
\href{https://www.youtube.com/watch?v=jBuvlWhNWkg}{\guillemotleft{}Homomorfismos inducidos\guillemotright{}}
(curso de Introducción a la Topología Algebraica). Una función continua que preserva puntos base induce un homomorfismo entre los grupos fundamentales, y esta asignación respeta composiciones e identidades.
\normalsize

\section*{Notación y enunciado}

Sea $f\colon X\to Y$ una función continua entre espacios topológicos, y sean $x_0\in X$ y $y_0\in Y$ puntos base con $f(x_0)=y_0$. Para ahorrar notación, en ese caso se escribe
\[
  f\colon (X,x_0)\to(Y,y_0),
\]
es decir, $f$ va de la pareja $X$ con su punto marcado $x_0$ a la pareja $Y$ con su punto marcado $y_0$.

\begin{teo}\label{teo:inducido}
Si $f\colon (X,x_0)\to(Y,y_0)$, entonces $f$ induce un homomorfismo
\[
  f_*\colon \pi_1(X,x_0)\to\pi_1(Y,y_0).
\]
Además:
\begin{enumerate}
  \item si $g\colon (Y,y_0)\to(Z,z_0)$ es continua, entonces $(g\circ f)_*=g_*\circ f_*$;
  \item $(\mathrm{id}_X)_*=\mathrm{id}_{\pi_1(X,x_0)}$.
\end{enumerate}
\end{teo}

\section*{La demostración}

\subsection*{La función inducida}

Primero se ve cómo se define la función inducida, que es fácil. Un elemento $[\gamma]\in\pi_1(X,x_0)$ es la clase de un lazo $\gamma\colon I\to X$ basado en $x_0$. Se puede componer con $f$: la composición $f\circ\gamma\colon I\to Y$ es un lazo basado en $y_0$ (pues $f(x_0)=y_0$). En el dibujo, el lazo $\gamma$ en $X$ se \emph{empuja} con $f$ y se obtiene el lazo $f\circ\gamma$ en $Y$: primero se aplica $\gamma$ y luego $f$.

\begin{center}
\begin{tikzpicture}[x=0.021cm,y=-0.021cm,line cap=round]
  \draw[thick] (130,110) ellipse (105 and 90);
  \draw[thick] (490,110) ellipse (105 and 90);
  \draw[very thick] (130,165) .. controls (50,165) and (55,55) .. (130,70) .. controls (205,55) and (210,165) .. (130,165);
  \draw[very thick,azul] (490,165) .. controls (410,155) and (425,50) .. (500,62) .. controls (570,80) and (565,175) .. (490,165);
  \draw[->,thick] (250,110) -- (370,110);
  \fill[extremo] (130,165) circle (4); \fill[extremo] (490,165) circle (4);
  \node at (130,188) {\small $x_0$};
  \node[extremo] at (490,188) {\small $y_0$};
  \node at (130,112) {\small $\gamma$};
  \node[azul] at (490,112) {\small $f\circ\gamma$};
  \node at (310,92) {\small $f$};
  \node at (30,30) {$X$}; \node at (590,30) {$Y$};
\end{tikzpicture}
\end{center}

\noindent Figura: el lazo $\gamma$ en $X$ y el lazo empujado $f\circ\gamma$ en $Y$.

Se define entonces
\[
  f_*\colon\pi_1(X,x_0)\to\pi_1(Y,y_0),\qquad f_*[\gamma]=[f\circ\gamma].
\]
Hay que verificar dos cosas: que $f_*$ está bien definida (pues $\pi_1$ es un conjunto de clases de equivalencia) y que es un homomorfismo de grupos.

\subsection*{$f_*$ está bien definida}

Sean $\gamma$ y $\gamma'$ lazos en $X$ basados en $x_0$ que son homotópicos relativos a los extremos; hay que ver que $f\circ\gamma$ y $f\circ\gamma'$ también lo son. Si $H\colon I\times I\to X$ es una homotopía relativa a los extremos entre $\gamma$ y $\gamma'$, lo más natural es considerar la composición
\[
  f\circ H\colon I\times I\to Y,
\]
y es un ejercicio muy sencillo verificar que es una homotopía relativa a los extremos entre los lazos empujados $f\circ\gamma$ y $f\circ\gamma'$. Esto demuestra que $f_*$ está bien definida: dos representantes de la misma clase en $\pi_1(X,x_0)$ dan la misma clase en $\pi_1(Y,y_0)$.

\subsection*{$f_*$ es un homomorfismo}

Hay que ver que $f_*[\gamma\cdot\beta]=f_*[\gamma]\cdot f_*[\beta]$, donde $\gamma\cdot\beta$ es el lazo que resulta de concatenar: primero se recorre $\gamma$ y luego $\beta$. Basta ver que empujar el concatenado es concatenar los empujados:
\[
  f\circ(\gamma\cdot\beta)=(f\circ\gamma)\cdot(f\circ\beta).
\]
En el dibujo es bastante claro: al empujar $\gamma$ con $f$ se obtiene $f\circ\gamma$, al empujar $\beta$ se obtiene $f\circ\beta$, y al empujar el concatenado se obtiene el concatenado de las imágenes de los lazos empujados. Lo único que restaría es escribir bien las parametrizaciones de la concatenación y de la imagen del concatenado, lo cual es un ejercicio muy sencillo.

\subsection*{La parte 2: la identidad}

Esta parte es muy sencilla. Un lazo $\gamma$ en $X$ empujado con la identidad es el mismo lazo, de modo que $(\mathrm{id}_X)_*[\gamma]=[\gamma]$ y no hay nada que demostrar.

\subsection*{La parte 1: composiciones}

Sean $f\colon(X,x_0)\to(Y,y_0)$ y $g\colon(Y,y_0)\to(Z,z_0)$, y la composición $g\circ f\colon (X,x_0)\to(Z,z_0)$. Sea $\gamma$ un lazo en $X$ basado en $x_0$. Para calcular $g_*\bigl(f_*[\gamma]\bigr)$ se empuja primero $\gamma$ con $f$ y luego el resultado con $g$, lo cual da el lazo $g\circ(f\circ\gamma)$. Como la composición de funciones es asociativa, esto es lo mismo que empujar $\gamma$ directamente con la composición:
\[
  g\circ(f\circ\gamma)=(g\circ f)\circ\gamma.
\]
Por lo tanto $(g\circ f)_*[\gamma]=g_*\bigl(f_*[\gamma]\bigr)$, es decir, $(g\circ f)_*=g_*\circ f_*$. Con esto termina la demostración del teorema. \qed

\section*{Lenguaje de categorías}

El teorema dice que las funciones inducidas se comportan a las mil maravillas: abren composiciones y la identidad va a la identidad. En el lenguaje de categorías, esto dice que el grupo fundamental es un \emph{funtor} de la categoría de espacios topológicos con puntos base a la categoría de grupos. Por ahora no se habla más de categorías; es solo una observación para quien sepa qué significa. Es un funtor de una categoría topológica a una categoría algebraica, y esto ya comienza a mostrar el espíritu de la topología algebraica: traducir problemas del mundo topológico al algebraico, es decir, al grupo fundamental. Lo que se hizo en los videos anteriores muestra que esto tiene aplicaciones directas a la topología y ayuda a entender más fenómenos topológicos.

\end{document}
